\documentclass[11pt]{article}
\usepackage{mystyle}

\title{Rosen --- \S 5.2: Strong Induction and Well-Ordering\\ \large Selected Exercises}
\author{Alston}
\date{Semester 2, 2026}

\begin{document}
\maketitle

% ======================================================================
\begin{exercise}{3}
    Let $P(n)$ be the statement that a postage of $n$ cents can be
    formed using just 3-cent stamps and 5-cent stamps. The parts of
    this exercise outline a strong induction proof that $P(n)$ is true
    for $n \geq 8$.
    \begin{enumerate}[label=(\alph*)]
        \item Show that the statements $P(8)$, $P(9)$, and $P(10)$ are
        true, completing the basis step of the proof.
        \item What is the inductive hypothesis of the proof?
        \item What do you need to prove in the inductive step?
        \item Complete the inductive step for $k \geq 10$.
        \item Explain why these steps show that this statement is true
        whenever $n \geq 8$.
    \end{enumerate}
\end{exercise}

% ======================================================================
\begin{exercise}{9}
    Use strong induction to prove that $\sqrt{2}$ is irrational.
    [\emph{Hint}: Let $P(n)$ be the statement that $\sqrt{2} \neq n/b$
    for any positive integer $b$.]
\end{exercise}

% ======================================================================
\begin{exercise}{11}
    Consider this variation of the game of Nim. The game begins with
    $n$ matches. Two players take turns removing matches, one, two, or
    three at a time. The player removing the last match loses. Use
    strong induction to show that if each player plays the best
    strategy possible, the first player wins if $n = 4j$, $4j+2$, or
    $4j+3$ for some nonnegative integer $j$ and the second player wins
    in the remaining case when $n = 4j+1$ for some nonnegative integer
    $j$.
\end{exercise}

% ======================================================================
\begin{exercise}{12}
    Use strong induction to show that every positive integer $n$ can
    be written as a sum of distinct powers of two, that is, as a sum
    of a subset of the integers $2^0 = 1$, $2^1 = 2$, $2^2 = 4$, and so
    on.
    [\emph{Hint}: For the inductive step, separately consider the case
    where $k+1$ is even and where it is odd. When it is even, note
    that $(k+1)/2$ is an integer.]
\end{exercise}

% ======================================================================
\begin{exercise}{13}
    A jigsaw puzzle is put together by successively joining pieces
    that fit together into blocks. A move is made each time a piece
    is added to a block, or when two blocks are joined. Use strong
    induction to prove that no matter how the moves are carried out,
    exactly $n - 1$ moves are required to assemble a puzzle with $n$
    pieces.
\end{exercise}

% ======================================================================
\begin{exercise}{15}
    Prove that the first player has a winning strategy for the game of
    Chomp, introduced in Example 12 in Section 1.8, if the initial
    board is square.
    [\emph{Hint}: Use strong induction to show that this strategy
    works. For the first move, the first player chomps all cookies
    except those in the left and top edges. On subsequent moves, after
    the second player has chomped cookies on either the top or left
    edge, the first player chomps cookies in the same relative
    positions in the left or top edge, respectively.]
\end{exercise}

% ======================================================================
\begin{exercise}{17}
    Use strong induction to show that if a simple polygon with at
    least four sides is triangulated, then at least two of the
    triangles in the triangulation have two sides that border the
    exterior of the polygon.
\end{exercise}

% ======================================================================
\begin{exercise}{19}
    Pick's theorem says that the area of a simple polygon $P$ in the
    plane with vertices that are all lattice points (that is, points
    with integer coordinates) equals $I(P) + B(P)/2 - 1$, where $I(P)$
    and $B(P)$ are the number of lattice points in the interior of $P$
    and on the boundary of $P$, respectively. Use strong induction on
    the number of vertices of $P$ to prove Pick's theorem.
    [\emph{Hint}: For the basis step, first prove the theorem for
    rectangles, then for right triangles, and finally for all
    triangles by noting that the area of a triangle is the area of a
    larger rectangle containing it with the areas of at most three
    triangles subtracted. For the inductive step, take advantage of
    Lemma 1.]
\end{exercise}

% ======================================================================
\begin{exercise}{24}
    \textit{Recall: a stable assignment pairs each suitor with a
    suitee so that no suitor and suitee who are not paired with each
    other both prefer each other to their current partners (as
    defined in the preamble to Exercise 60 in Section 3.1).}

    A stable assignment is called \emph{optimal for suitors} if no
    stable assignment exists in which a suitor is paired with a
    suitee whom this suitor prefers to the person to whom this suitor
    is paired in this stable assignment. Use strong induction to show
    that the deferred acceptance algorithm produces a stable
    assignment that is optimal for suitors.
\end{exercise}

% ======================================================================
\begin{exercise}{27}
    Show that if the statement $P(n)$ is true for infinitely many
    positive integers $n$ and $P(n+1) \to P(n)$ is true for all
    positive integers $n$, then $P(n)$ is true for all positive
    integers $n$.
\end{exercise}

% ======================================================================
\begin{exercise}{36}
    The well-ordering property can be used to show that there is a
    unique greatest common divisor of two positive integers. Let $a$
    and $b$ be positive integers, and let $S$ be the set of positive
    integers of the form $as + bt$, where $s$ and $t$ are integers.
    \begin{enumerate}[label=(\alph*)]
        \item Show that $S$ is nonempty.
        \item Use the well-ordering property to show that $S$ has a
        smallest element $c$.
        \item Show that if $d$ is a common divisor of $a$ and $b$,
        then $d$ is a divisor of $c$.
        \item Show that $c \mid a$ and $c \mid b$.
        [\emph{Hint}: First, assume that $c \nmid a$. Then
        $a = qc + r$, where $0 < r < c$. Show that $r \in S$,
        contradicting the choice of $c$.]
        \item Conclude from (c) and (d) that the greatest common
        divisor of $a$ and $b$ exists. Finish the proof by showing
        that this greatest common divisor is unique.
    \end{enumerate}
\end{exercise}

% ======================================================================
\begin{exercise}{39}
    Can you use the well-ordering property to prove the statement:
    ``Every positive integer can be described using no more than
    fifteen English words''? Assume the words come from a particular
    dictionary of English.
    [\emph{Hint}: Suppose that there are positive integers that cannot
    be described using no more than fifteen English words. By
    well-ordering, the smallest positive integer that cannot be
    described using no more than fifteen English words would then
    exist.]
\end{exercise}

% ======================================================================
\begin{exercise}{40}
    Use the well-ordering property to show that if $x$ and $y$ are
    real numbers with $x < y$, then there is a rational number $r$
    with $x < r < y$.
    [\emph{Hint}: Use the Archimedean property to find a positive
    integer $A$ with $A > 1/(y-x)$. Then show that there is a rational
    number $r$ with denominator $A$ strictly between $x$ and $y$.]
\end{exercise}

% ======================================================================
\begin{exercise}{42}
    Show that the principle of mathematical induction and strong
    induction are equivalent; that is, each can be shown to be valid
    from the other.
\end{exercise}

\end{document}